Teacher-based assessments dictated university admissions in the UK during the COVID 19 pandemic, instead of traditional standardized exams (\emph{A-levels}). The British governments announced the switch after universities and students had already agreed on the grade requirements for admission, as the British university replies to students applications with legally binding contracts that state the bar of grades for admission. The government resorted to teachers to provide a replacement grade for students A-level outcome. As a consequence, the policy led to a significant degree of grade inflation. 

This study explores the patterns of grade inflation in various types of schools and among different demographics of students. British private schools, known as \emph{Independent Schools}, rely mainly on tuition from wealthy families, while public schools are funded through allocations from local governments and cater to local students. Allegations arose against private secondary schools for inflating grades to enhance their pupils' placement statistics\footnote{Refer to "A-levels 2021: Rise three times more in private schools," available at \url{https://www.tes.com/magazine/news/secondary/levels-2021-rise-three-times-more-private-schools}. Accessed on December 18, 2024.} While improvements in female students' grades compared to their male counterparts were also noted\footnote{See "Boys set to close gap with girls for top A-level grades," available at \url{https://www.thetimes.com/uk/education/article/a-level-results-2024-boys-expected-higher-girls-kc7sqkpqt}. Accessed on December 29, 2024.}. Moreover, A-levels provide students with a broad selection of subjects, all assessed centrally. Universities tend to favor some subjects over others, emphasizing the trustworthiness of exam outcomes. The subjects traditionally deemed essential by universities, such as mathematics and English literature, are referred to as \emph{Facilitating Subjects}, in contrast to \emph{Non-Facilitating subjects}, which includes subjects such as drama studies.

I estimate the increments in probability in obtaining a top grade using an interrupted time series design with administrative data on university application. The slackened grade standards reportedly elevated grades across students with vast background. However, the impacts for affluent students is not easily estimated, since students differ in their probabilities of obtaining top grades. Therefore, I account for differences in student backgrounds by comparing the probability of obtaining a top grade across different grading methods, rather than directly estimating the changes in the probability. Moreover, I include rich demographic details and performance in previous standardized exams (\emph{GCSE}) in the estimation to control for the diverse backgrounds of students. Shifting focus, the sudden timing of the pandemic ensures that my estimates are driven from the method change rather than shifts in student composition across the pre/post-pandemic years. Moreover, since all university applicants take GCSE exams, my estimates are comparable across students attending different schools. 

I find differential increments of final grades across school types and their pupil quality. I find that independent schools inflated grades by 0.1 standard deviation points compared to other public schools. I find an inverse U shape pattern between the share of students within a school with good scores in the previous national exams and the extent of their inflation. I observe a negative relationship between schools, with less than half of their students scoring good grades on their national exams. A school in the 50th percentile of the quality distribution assigned 0.5 standard deviation points higher grades than the school in the 0th percentile. A positive relationship is observed for schools with more than half of its students performing well. Additionally, I find grade inflation occurred more often for students taking non-quantitative subjects. Exams in facilitating subject categories were less inflated than non-facilitating subjects. Non-facilitating subjects had 9 percentage points less probability of obtaining an A or $A^*$ than facilitating subjects. However, facilitating subjects' exams were also inflated with their probability of receiving an A or $A^*$ increasing by 8 percentage points compared to previous years. I do not find differences in increments of assigned grades across student demographics after controlling for the student's subject choice.

My findings highlight how privately funded schools, such as independent schools, react to opportunities to improve the grades
of their students faster than other schools. In contrast to public schools, independent schools run by privately raised tuition paid by the parents of their students. Independent schools compete with each other by the number of students they place in better universities. Infer that independent schools took advantage of the loophole in grade assignment more quickly than other types of schools. 

My results also demonstrate how teacher-based grading rewards students who take subjects with more subjective discretion over their grading. Teachers may find difficulties in ranking performances of subjective subjects, due to their own preferences of content and the non-numerical responses of students. Teachers may also be reluctant to assign bad grades to students when they cannot strictly rank students from each other. The smaller increments in grades for quantitative subjects indicate that when teachers are aware of the students' ability with a well-defined metric, they are less likely to deviate from the grades a student would receive in a standardized exam. 



My results imply that the beneficiaries of non-standardized exam spreads across a diverse range of student bodies. Students in underrepresented demographic groups benefit from a teacher-based grading scheme since such subjects are difficult for teachers to assign bad grades. Their application success rate should improve because many underrepresented students take non-facilitating subjects. I infer students from advantageous backgrounds will not switch to non-facilitating subjects to increase their average grade, since many elite universities have strict subject requirements. There should be a reshuffling of university assignments between students from advantageous backgrounds who did not perform so well in their A-level with ambitious subject choice and students from less affluent backgrounds with less ambitious subject choice but with better overall grades. However, teacher-based grading is prone to manipulation, by more affluent schools having a stronger incentive to inflate the pupils' grades. 
\clearpage